\section{Évaluation à l’échelle de Pell et fonctionnelle constitutive de Barbero--Immirzi}
\label{iii:sec:pell-barbero}

La relation entre la réponse du vide et le moment de Catalan permet de
composer deux niveaux du transport : la profondeur pondérée du caractère
de quart de tour et l’évaluation des canaux angulaires. L’échelle de
Pell distingue un argument de la famille de moments ; les canaux du
vide reçoivent, en revanche, les arguments déterminés par la paire angulaire.
La fonctionnelle de Barbero--Immirzi s’exprime sur ces derniers au moyen de
la même réponse constitutive et de son inverse. Le lien s’établit ainsi
entre des opérations définies, en conservant la position de chaque évaluation
dans la construction APP--TRIT--TPK.

\subsection{Moment orienté et évaluation dans l’unité de Pell}

Écrivons
\begin{equation}
 G_{\rm Cat}:=\mathcal C_2(1)
 =\sum_{k\geq0}\frac{(-1)^k}{(2k+1)^2},\qquad
 \lambda_{\rm P}=2+\sqrt3,\qquad
 \rho_{\rm P}=2-\sqrt3=\lambda_{\rm P}^{-1}.
 \label{iii:eq:pell-scales}
\end{equation}
La constante $G_{\rm Cat}$ est la valeur à l’extrémité du moment construit dans
\eqref{iii:eq:c2}. La paire $\lambda_{\rm P},\rho_{\rm P}$ correspond aux
orientations expansive et contractante de l’unité de Pell :
$(2+\sqrt3)(2-\sqrt3)=1$. L’identité de norme $2^2-3\cdot1^2=1$
fixe la réciprocité. La spécialisation du moment en $\rho_{\rm P}$
est postérieure à la construction de son caractère orienté.

L’équation différentielle \eqref{iii:eq:c2-integral}, avec condition nulle
à l’origine, donne
\begin{equation}
 \mathcal D\mathcal C_2(s)=\arctan s,\qquad
 \mathcal C_2(s)=\int_0^s\frac{\arctan t}{t}\,dt,
 \qquad \mathcal D=s\frac{d}{ds}.
 \label{iii:eq:pell-moment-integral}
\end{equation}
En effet, intégrer $\mathcal D^2\mathcal C_2=s/(1+s^2)$ par rapport à
$ds/s$ produit $\arctan s$ ; une seconde intégration produit l’expression
indiquée. L’intégrande se prolonge continûment en $t=0$ avec la valeur $1$.
La forme angulaire de cette intégrale permet de déterminer exactement
l’évaluation en $\rho_{\rm P}$.

\begin{lemma}[Intégrale angulaire du caractère impair]
\label{iii:lem:log-tangent}
Pour $0<\theta\leq\pi/4$,
\begin{equation}
 -\int_0^\theta\log(\tan x)\,dx
 =\sum_{\substack{n\geq1\\n\text{ impair}}}
      \frac{\sin(2n\theta)}{n^2}.
 \label{iii:eq:log-tangent}
\end{equation}
\end{lemma}
\begin{proof}
Pour $0<a<1$, le développement du logarithme à l’intérieur du disque unité implique
\begin{align*}
 L_a(x)
 &:=\log|1-ae^{2ix}|-\log|1+ae^{2ix}|\\
 &=-2\sum_{\substack{n\geq1\\n\text{ impair}}}
       \frac{a^n\cos(2nx)}{n}.
\end{align*}
La série converge uniformément sur $x$ pour chaque $a<1$ ; son intégration
terme à terme fournit
\[
 -\int_0^\theta L_a(x)\,dx
 =\sum_{\substack{n\geq1\\n\text{ impair}}}
       \frac{a^n\sin(2n\theta)}{n^2}.
\]
Lorsque $a\uparrow1$, le quotient des modules converge vers
$|1-e^{2ix}|/|1+e^{2ix}|=\tan x$ pour $0<x\leq\theta$.
La singularité du logarithme à l’origine est intégrable. Pour justifier
la limite, on peut restreindre $a\geq1/2$ et observer que
\[
 |1-ae^{2ix}|^2=(1-a)^2+4a\sin^2x
 \geq2\sin^2x,
 \qquad 1\leq|1+ae^{2ix}|\leq2
\]
sur $0\leq x\leq\pi/4$. Avec la borne supérieure
$|1-ae^{2ix}|\leq2$, ces inégalités donnent
$|L_a(x)|\leq c+|\log x|$ avec une constante indépendante de $a$.
La convergence dominée permet de passer à la limite dans l’intégrale.
Dans la série intégrée, on applique la domination par $\sum n^{-2}$.
Les deux limites prouvent \eqref{iii:eq:log-tangent}.
\end{proof}

\begin{theorem}[Évaluation Catalan--Pell]
\label{iii:thm:catalan-pell}
Le moment orienté satisfait
\begin{align}
 \mathcal C_2(\rho_{\rm P})
 &=\frac23G_{\rm Cat}-\frac{\pi}{12}\log\lambda_{\rm P},
 \label{iii:eq:catalan-pell}\\
 \mathcal D\mathcal C_2(\rho_{\rm P})&=\frac{\pi}{12},&
 \mathcal D^2\mathcal C_2(\rho_{\rm P})&=\frac14.
 \label{iii:eq:catalan-pell-derivatives}
\end{align}
\end{theorem}
\begin{proof}
La substitution $t=\tan x$ dans \eqref{iii:eq:pell-moment-integral},
suivie d’une intégration par parties, donne
\begin{equation}
 \mathcal C_2(\tan\theta)
 =\theta\log(\tan\theta)
   -\int_0^\theta\log(\tan x)\,dx.
 \label{iii:eq:angular-moment}
\end{equation}
Le terme de frontière en zéro s’annule parce que $x\log(\tan x)\to0$.
La formule de la tangente d’une différence fournit
\[
 \tan\frac{\pi}{12}
 =\tan\left(\frac\pi4-\frac\pi6\right)
 =\frac{\sqrt3-1}{\sqrt3+1}=2-\sqrt3.
\]
Il reste à évaluer la série de \eqref{iii:eq:log-tangent} en
$\theta=\pi/12$. Soit $\chi_{-4}(n)=(-1)^{(n-1)/2}$ pour $n$ impair.
Dans les six classes impaires modulo $12$, on vérifie
\begin{equation}
 2\sin\frac{n\pi}{6}
 =\chi_{-4}(n)+3\,\mathbf1_{3\mid n}\chi_{-4}(n/3).
 \label{iii:eq:character-pell}
\end{equation}
La vérification finie est
\[
\begin{array}{c|rrrrrr}
 n\bmod12&1&3&5&7&9&11\\ \hline
 2\sin(n\pi/6)&1&2&1&-1&-2&-1\\
 \chi_{-4}(n)&1&-1&1&-1&1&-1\\
 3\mathbf1_{3\mid n}\chi_{-4}(n/3)&0&3&0&0&-3&0
\end{array}
\]
et les deux membres sont périodiques modulo $12$ sur les indices impairs.
La convergence absolue permet de séparer les termes et d’introduire $n=3m$
dans le second terme :
\begin{align*}
 \sum_{\substack{n\geq1\\n\text{ impair}}}
       \frac{\sin(n\pi/6)}{n^2}
 &=\frac12\sum_{n\text{ impair}}\frac{\chi_{-4}(n)}{n^2}
  +\frac32\sum_{m\text{ impair}}\frac{\chi_{-4}(m)}{(3m)^2}\\
 &=\left(\frac12+\frac16\right)G_{\rm Cat}
 =\frac23G_{\rm Cat}.
\end{align*}
Les sommes sans borne inférieure dans cette expression parcourent les entiers
positifs. En substituant dans \eqref{iii:eq:angular-moment} et en utilisant
$\log\rho_{\rm P}=-\log\lambda_{\rm P}$, on obtient
\eqref{iii:eq:catalan-pell}. La première dérivée résulte de
$\arctan\rho_{\rm P}=\pi/12$ ; la seconde, de
\[
 \frac{\rho_{\rm P}}{1+\rho_{\rm P}^2}
 =\frac1{\rho_{\rm P}+\rho_{\rm P}^{-1}}=\frac14.
\]
\end{proof}

Le coefficient $2/3$ provient de la décomposition du caractère dans
\eqref{iii:eq:character-pell} : la restriction aux multiples de $3$
modifie le poids de profondeur par le facteur $3^{-2}$. L’évaluation
réunit donc le quart de tour et la subdivision dodécaphasique dans une
identité exacte de moments. Les coefficients sont fixés avant toute
comparaison décimale.

\subsection{Composition de phase et profondeur pondérée}

La représentation de profondeur précise ce que conserve cette évaluation.
Dans $\ell^2(\mathbb N_{>0})$, définissons
\[
 N\delta_n=n\delta_n,\qquad U_4\delta_n=i^n\delta_n,
 \qquad Q_4=\frac{U_4-U_4^*}{2i}.
\]
Pour $0<s<1$, l’opérateur $N^{-2}s^NQ_4$ est à trace, et
\begin{equation}
 \mathcal C_2(s)=\operatorname{Tr}(N^{-2}s^NQ_4).
 \label{iii:eq:pell-depth-trace}
\end{equation}
Cette égalité s’obtient en évaluant sur la base $\delta_n$ ; les
entrées diagonales sont $s^n\sin(n\pi/2)/n^2$. L’échange
$U_4\leftrightarrow U_4^*$ inverse $Q_4$ et, avec lui, le moment
orienté. L’indice $n$ et la pondération $s^n$ restent fixes.

La composition s’effectue avant de prendre la trace. Pour
$a\in\mathbb Z/4\mathbb Z$, soit $B(s,a)=s^NU_4^a$. Comme les deux
facteurs sont diagonaux dans la même base,
\begin{equation}
 B(s,a)B(t,b)=B(st,a+b),\qquad
 B(s,1)^4=s^{4N}.
 \label{iii:eq:pell-depth-composition}
\end{equation}
Le retour de phase conserve la contraction accumulée. En particulier,
$B(\rho_{\rm P},1)^4=\rho_{\rm P}^{4N}$, dont la valeur propre à la profondeur
$n$ est $\rho_{\rm P}^{4n}$. Cette réalisation donne un sens
opératoriel à l’évaluation contractante ; la somme scalaire du moment
agrège les profondeurs après avoir conservé leur indice et leur orientation.
Les arguments constitutifs $s_\pm=q_\pm^{30}$ restent déterminés par
la paire angulaire de \eqref{iii:eq:recover-angular}. Les formules précédentes
ne contiennent pas d’égalité entre ces arguments et $\rho_{\rm P}$.

\subsection{Incidence dodécaphasique et fonctionnelle de retour}

La contribution de Barbero--Immirzi utilise le transport angulaire et
l’incidence hexade--octade. Soient $\mathsf H\subset\mathsf O$ une hexade
et une octade avec la paire marquée de l’hexade conservée. Les ensembles de
drapeaux pertinents sont
\begin{equation}
 \mathcal F_6=\mathsf H\times\binom{\mathsf H}{2}
 \lhook\joinrel\longrightarrow
 \mathcal F_8=\mathsf O\times\binom{\mathsf H}{2},
 \qquad |\mathcal F_6|=90,\quad |\mathcal F_8|=120.
 \label{iii:eq:barbero-flags}
\end{equation}
Les cardinaux sont $6\binom62$ et $8\binom62$ ; la même incidence
conserve le défaut normalisé
$(90-120)/(24\binom62)=-1/12$. Ce dénombrement identifie les deux secteurs
employés par le retour, tandis que le calendrier détermine leurs
multiplicités relatives.

Pour $m\in\{90,120\}$, le retour tritique sur un canal contractant est
\begin{equation}
 S_m(q)=\sum_{j\geq0}q^{m+3mj}
       =\frac{q^m}{1-q^{3m}},\qquad 0<q<1.
 \label{iii:eq:barbero-return-series}
\end{equation}
Chaque terme correspond à la hauteur initiale $m$ et à $j$ retours de
longueur $3m$. La série géométrique conserve donc tant la hauteur
initiale que la période du retour. Le calendrier dodécaphasique
$\mathbb Z/12\mathbb Z$ apporte douze secteurs et une unique incidence
orientée de frontière, notée $\partial_{\rm ret}$. Dans le module
libre des événements, sa chaîne fondamentale est
\begin{equation}
 c_{12}^{+}
 =\sum_{j\in\mathbb Z/12\mathbb Z}[j,\mathcal F_6]
  -[\partial_{\rm ret},\mathcal F_8].
 \label{iii:eq:barbero-cycle}
\end{equation}
Le lecteur linéaire $\mathscr L_q$ assigne $S_{90}(q)$ à chaque générateur
$[j,\mathcal F_6]$ et $S_{120}(q)$ au générateur de frontière. On obtient
\begin{equation}
 \mathscr L_q(c_{12}^{+})=12S_{90}(q)-S_{120}(q).
 \label{iii:eq:barbero-weights}
\end{equation}
Les poids $12$ et $-1$ sont les multiplicités orientées de la chaîne.
L’orientation conjuguée satisfait $c_{12}^{-}=-c_{12}^{+}$, avec les
mêmes types d’incidence.

\begin{proposition}[Coefficient singulier du retour fondamental]
\label{iii:prop:barbero-pole}
Le coefficient de $(1-q)^{-1}$ de
$\mathscr L_q(c_{12}^{+})$ est $1/24$.
\end{proposition}
\begin{proof}
Pour chaque entier $m>0$, la factorisation
$1-q^{3m}=(1-q)\sum_{j=0}^{3m-1}q^j$ donne
\[
 \lim_{q\uparrow1}(1-q)S_m(q)=\frac1{3m}.
\]
Appliquer cette identité à la fonctionnelle déjà déterminée dans
\eqref{iii:eq:barbero-weights} produit
\begin{equation}
 \lim_{q\uparrow1}(1-q)\mathscr L_q(c_{12}^{+})
 =\frac{12}{270}-\frac1{360}=\frac1{24}.
 \label{iii:eq:barbero-pole}
\end{equation}
Dans la coordonnée complexe $q-1$, le résidu est $-1/24$.
\end{proof}

L’évaluation singulière vérifie une conséquence de l’incidence et du
retour tritique. L’ordre démonstratif est fixé par
\eqref{iii:eq:barbero-flags}--\eqref{iii:eq:barbero-weights} ; le
coefficient singulier se calcule après la construction de la chaîne.

\subsection{Factorisation à travers l’opérateur constitutif}

Conservons la coordinatisation angulaire relevée de la
section~\ref{iii:sec:hojas}. Pour rendre explicite la conversion angulaire,
écrivons
\begin{equation}
 A_{\deg}=\frac{180x}{\pi},\qquad
 C^*_{\deg}=\frac{180y}{\pi},\qquad
 q_\pm=e^{-x\mp y},\qquad s_\pm=q_\pm^{30}.
 \label{iii:eq:barbero-angular-lift}
\end{equation}
L’orientation considérée dans \eqref{iii:eq:rchannels} correspond à
$0<y<x$ ; l’échange de feuillets s’exprime par $y\mapsto-y$ sur le
domaine plus large $x>|y|$. Les exponentielles réelles agissent sur ces
relèvements, qui conservent les valeurs de $x\pm y$ pendant la
composition.

La fonctionnelle orientée de Barbero--Immirzi est
\begin{equation}
 \gamma=\frac{\pi A_{\deg}C^*_{\deg}}{180}
   +12\bigl[S_{90}(q_-)-S_{90}(q_+)\bigr]
   -\bigl[S_{120}(q_-)-S_{120}(q_+)\bigr].
 \label{iii:eq:barbero-functional}
\end{equation}
La première composante est l’évaluation angulaire ; la seconde est la
différence du lecteur de la chaîne fondamentale dans les canaux conjugués.
Les arguments sont reçus du transport angulaire précédemment construit.
La formule détermine un scalaire adimensionnel sur cette structure de
départ.

Soit $\psi:(3/4,1)\to(0,1)$ l’inverse de $f$ construit au moyen de
\eqref{iii:eq:cubic}. Définissons
\begin{equation}
 \mathcal H(s)=\frac{12s^3}{1-s^9}
             -\frac{s^4}{1-s^{12}}
             -\frac{(\log s)^2}{20\pi},\qquad 0<s<1.
 \label{iii:eq:barbero-h}
\end{equation}
Cette fonction est analytique réelle sur son domaine. Les exposants $3,4,9,12$
résultent de l’expression des hauteurs $90,120$ et des retours $270,360$ en
unités de $30$ ; le coefficient de la partie logarithmique est fixé lorsqu’on
transporte la contribution angulaire.

\begin{theorem}[Expression constitutive orientée de Barbero--Immirzi]
\label{iii:thm:barbero-factorization}
Dans la coordinatisation \eqref{iii:eq:barbero-angular-lift},
\begin{equation}
 \gamma=\mathcal H(s_-)-\mathcal H(s_+)
       =\mathcal H\bigl(\psi(r_-)\bigr)
        -\mathcal H\bigl(\psi(r_+)\bigr).
 \label{iii:eq:barbero-factorization}
\end{equation}
Sur la représentation bidimensionnelle des canaux,
\begin{equation}
 \gamma=-\operatorname{Tr}
    \left[\mathcal R\,\mathcal H\bigl(\psi(\mathcal V)\bigr)\right],
 \label{iii:eq:barbero-trace}
\end{equation}
où $\operatorname{Tr}$ est la trace ordinaire, avec
$\operatorname{Tr}I=2$.
\end{theorem}
\begin{proof}
La substitution $s=q^{30}$ dans \eqref{iii:eq:barbero-return-series} donne
\[
 S_{90}(q)=\frac{s^3}{1-s^9},\qquad
 S_{120}(q)=\frac{s^4}{1-s^{12}}.
\]
De plus, \eqref{iii:eq:barbero-angular-lift} implique
\[
 \log s_\pm=-30(x\pm y)
 =-\frac\pi6(A_{\deg}\pm C^*_{\deg}).
\]
Par différence de carrés,
\begin{align}
 \frac{(\log s_+)^2-(\log s_-)^2}{20\pi}
 &=\frac{900\bigl[(x+y)^2-(x-y)^2\bigr]}{20\pi}\\
 &=\frac{180xy}{\pi}
 =\frac{\pi A_{\deg}C^*_{\deg}}{180}.
 \label{iii:eq:barbero-log-term}
\end{align}
La somme de ces identités prouve la première égalité de
\eqref{iii:eq:barbero-factorization}. La seconde utilise l’inversion
$\psi(f(s_\pm))=s_\pm$ démontrée dans la
section~\ref{iii:sec:constitutiva}.

Pour l’expression opératorielle, les projecteurs de feuillet sont de rang un et
$\mathcal RP_\pm=\pm P_\pm$. Le calcul spectral fournit
\[
 \mathcal H\bigl(\psi(\mathcal V)\bigr)
 =\mathcal H(s_+)P_++\mathcal H(s_-)P_-.
\]
Multiplier par $\mathcal R$ et prendre la trace ordinaire donne
$\mathcal H(s_+)-\mathcal H(s_-)$. Son opposé est $\gamma$.
\end{proof}

La marque d’orientation $\mathcal R$ fait partie de la fonctionnelle. Si
l’on maintient cette marque fixe et que l’on échange les canaux de $\mathcal V$,
le signe de $\gamma$ change. L’évaluation du spectre sans ses étiquettes conserve
la paire non ordonnée, tandis que \eqref{iii:eq:barbero-trace} conserve la
contribution orientée. Si l’on emploie la trace normalisée
$\tau_2=\operatorname{Tr}/2$, la même formule s’écrit
$\gamma=-2\tau_2[\mathcal R\mathcal H(\psi(\mathcal V))]$.

En revanche, un changement simultané de représentation
$(\mathcal R,\mathcal V)\mapsto
(U\mathcal RU^{-1},U\mathcal VU^{-1})$ conserve la fonctionnelle :
le calcul spectral se transporte par conjugaison et la trace est
invariante. Échanger la réponse par rapport à une orientation fixe
et transporter conjointement la réponse et l’orientation sont donc
des opérations différentes.

Cette normalisation est compatible avec les amplifications de
\eqref{iii:eq:logvac}. Pour
$\mathcal V_m=\mathcal V\otimes I_{9^m}$ et
$\mathcal R_m=\mathcal R\otimes I_{9^m}$,
\begin{equation}
 \gamma=-\frac1{9^m}\operatorname{Tr}
 \left[\mathcal R_m\mathcal H\bigl(\psi(\mathcal V_m)\bigr)\right].
 \label{iii:eq:barbero-amplification}
\end{equation}
En effet, le calcul fonctionnel conserve le produit tensoriel et la trace
de $I_{9^m}$ est $9^m$. La multiplicité de la représentation se sépare
de l’évaluation orientée.

\subsection{Reconstruction à partir des coordonnées du vide}

Les identités \eqref{iii:eq:four-identities} permettent d’exprimer la même
fonctionnelle au moyen de l’impédance et de la vitesse normalisées :
\begin{align}
 \gamma={}&
 \mathcal H\!\left(\psi\!\left(
     \sqrt{\frac{\widehat Z}{\widehat c}}\right)\right)
 -\mathcal H\!\left(\psi\!\left(
     \frac1{\sqrt{\widehat Z\widehat c}}\right)\right).
 \label{iii:eq:barbero-zc}
\end{align}
Le domaine positif de cette reconstruction est décrit par
\begin{equation}
 1<\widehat Z\widehat c<\frac{16}{9},\qquad
 \frac9{16}<\frac{\widehat Z}{\widehat c}<1.
 \label{iii:eq:barbero-zc-domain}
\end{equation}
En effet, ces inégalités équivalent à
$3/4<r_\pm<1$ lorsqu’on substitue les racines positives de
\eqref{iii:eq:four-identities}. Pour l’orientation $y>0$, on conserve
en outre $r_-<r_+$, ou de manière équivalente $\widehat Z<1$. Les points du
domaine sont utilisés avec la compatibilité angulaire déjà établie par
le TPK ; l’inversion est une reconstruction ultérieure de ses canaux.
Les coordonnées sont les coordonnées internes de \eqref{iii:eq:four} : un changement
d’unités doit être défait au moyen de \eqref{iii:eq:undo-units} avant
d’appliquer $\psi$.

L’échange $y\mapsto-y$ agit comme
$(\widehat Z,\widehat c)\mapsto(\widehat Z^{-1},\widehat c)$ et
produit $\gamma\mapsto-\gamma$. Dans le cas équilibré $y=0$,
les deux canaux coïncident et \eqref{iii:eq:barbero-factorization} donne
$\gamma=0$. La vitesse normalisée conserve le produit des
canaux ; sa lecture conjointe avec l’impédance retrouve les deux
réponses et l’orientation nécessaire à l’évaluation de la fonctionnelle.
Dans ce cas équilibré, $\mathcal V=rI$ a un spectre dégénéré :
la marque $\mathcal R$ est conservée comme donnée de représentation, bien que
le scalaire $r$ ne distingue plus ses projecteurs. Pour un spectre simple,
en revanche,
$P_+=(\mathcal V-r_-I)/(r_+-r_-)$ et $P_-=I-P_+$.

Le résultat central de cette composition est la factorisation
\eqref{iii:eq:barbero-factorization}. Le moment de Catalan apporte le
coefficient orienté qui détermine $f$ ; l’inversion cubique retrouve
l’argument de chaque canal ; l’incidence dodécaphasique et la paire angulaire
déterminent $\mathcal H$. L’échelle de Pell distingue une évaluation
exacte de la famille, avec la preuve de
\eqref{iii:eq:catalan-pell}. Chaque relation conserve son argument et son
opération : la constante à l’extrémité, le moment contractant et la fonctionnelle
constitutive constituent des évaluations reliées de structures
spécifiées.
